\documentclass[../../main.tex]{subfiles}
\begin{document}
\section{简单解法}\label{section:appD-simple-solution-techniques}

为了从本书中获得最大收益，你需要能够求出少数几个微分方程与微分方程组的显式解。你大概已学过这类方程的多种解法。下列简单类型的方程对本书的需要已绰绰有余。

\subsection*{可分离变量方程}

单个函数 $y(t)$ 的一阶微分方程，如果能写成形式
\begin{align*}
y'(t)=f(y(t))g(t),
\end{align*}
其中 $f,g$ 连续且 $f$ 不恒为零，就称为\emph{可分离变量}的。任意可分离方程都可（至少原则上）通过除以 $f(y(t))$、两边积分、并用变量代换变换左端积分来求解：
\begin{align*}
\frac{y'(t)}{f(y(t))}=g(t),\qquad
\int\frac{y'(t)\,\mathrm{d}t}{f(y(t))}=\int g(t)\,\mathrm{d}t,\qquad
\int\frac{\mathrm{d}y}{f(y)}=\int g(t)\,\mathrm{d}t.
\end{align*}
如果所得不定积分能显式算出，结果就是一个涉及 $y$ 与 $t$ 的关系式，可（再次，原则上）对 $y$ 解出。积分常数随后可调整以达到 $y$ 所需的初值。可分离方程包括形如 $y'(t)=g(t)$ 的方程作为平凡特例，它可通过直接积分解出。

\subsection*{$2\times2$ 常系数线性方程组}

形如
\begin{align}\label{align:appD-linear-system}
x'(t)=ax(t)+by(t),\qquad y'(t)=cx(t)+dy(t),
\end{align}
（其中 $a,b,c,d$ 是常数）的方程组，可用矩阵记号写成 $Z'(t)=AZ(t)$，其中
\begin{align*}
Z(t)=\begin{pmatrix}x(t)\\ y(t)\end{pmatrix},\qquad
A=\begin{pmatrix}a&b\\ c&d\end{pmatrix}.
\end{align*}
这种方程组的解集总构成 $\mathbb R$ 上的一个二维向量空间。一旦找到两个线性无关解，每个其他解都是它们的线性组合，常数可调整以匹配任意初值。

总可能找到至少一个（或许复值）形如 $Z(t)=e^{\lambda t}Z_0$ 的解，其中 $\lambda$ 是 $A$ 的特征值、$Z_0$ 是相应特征向量。如果 $A$ 有两个不同特征值，则有两个这样的解，它们张成解空间。（如果初值为实，则相应解为实。）另一方面，如果 $A$ 只有一个特征值，则有两种情形。如果 $A-\lambda I_2\neq0$，则存在向量 $Z_1$（称为广义特征向量）使 $(A-\lambda I_2)Z_1=Z_0$，第二个线性无关解由 $Z(t)=e^{\lambda t}(tZ_0+Z_1)$ 给出。否则 $A=\lambda I_2$，\eqref{align:appD-linear-system} 中两方程解耦，可独立求解。

\subsection*{部分解耦方程组}

如果\eqref{align:appD-system} 中某一个微分方程（例如 $\dot y^i(t)$ 的方程）除 $y^i(t)$ 外不涉及任何其他因变量，则可先尝试求解该方程，再把解代入其他方程，从而得到一个未知函数更少的方程组，它也许可用上述某一方法求解。

\begin{exercise}\label{exercise:appD-D-7}
求解下列初值问题。
\begin{enumerate}[(a)]
\item $x'(t)=x(t)^2$；$x(0)=x_0$。
\item $x'(t)=\dfrac{1}{x(t)}$；$x(0)=x_0>0$。
\item $x'(t)=y$；$x(0)=x_0$；$y'(t)=1$；$y(0)=y_0$。
\item $x'(t)=x$；$x(0)=x_0$；$y'(t)=2y$；$y(0)=y_0$。
\item $x'(t)=-y$；$x(0)=x_0$；$y'(t)=x$；$y(0)=y_0$。
\item $x'(t)=-x(t)+y(t)$；$x(0)=x_0$；$y'(t)=-x(t)-y(t)$；$y(0)=y_0$。
\item $x'(t)=1$；$x(0)=x_0$；$y'(t)=\dfrac{1}{1+x(t)^2}$；$y(0)=y_0$。
\end{enumerate}
\end{exercise}
\end{document}
